\section{Verifying: Two-Property Translation Validation for a Real Pipeline}

% OUTLINE ONLY. Source of truth: prototypes/fcvd_ct,
% docs/research/polygeist-verification.agents.md (the loop journal),
% docs/research/fcvd-selfcomposition.agents.md,
% docs/research/fcvd-verification-stages.agents.md.

\begin{itemize}
  \item Why Polygeist: most translatable corpus of the six compilers mapped
    (61.9\% of operation mentions vs 53.4\% HEIR, 17.3\% torch-mlir); speaks
    the shared dialects (affine, scf, memref, arith) --- the compiler most
    likely to be verified end-to-end rather than in patches.
  \item The property: a lowering is VERIFIED only if \emph{both halves} hold ---
    semantic congruence (computes the same thing) and non-interference
    (self-composition; hole congruence includes poison bits).
  \item The method: one checked specification per pass, transcribed from the
    pass source with a citation, each preserving verdict paired with a
    falsifying control; append-only journal; coverage counter re-checks every
    verdict.
  \item Pipeline status table (regenerate from live data at submission):
    eight steps from \texttt{tools/cgeist/driver.cc} --- mem2reg,
    loop-restructure, affine-cfg, canonicalize-for, lower-affine,
    parallel-lower, scf-to-openmp, polygeist-to-llvm; note which are checked
    vs model-blocked (the two parallel steps).
  \item Scaling the check: the taint prefilter before the solver; SMT semantics
    for LLVM memory ops closing the memref$\rightarrow$LLVM slice.
  \item Per-program vs per-pass proofs: macro-templates cover a translation
    only if both halves hold; \texttt{affine.load}/\texttt{affine.store} with
    identity maps as the cheapest-and-most-used translations.
\end{itemize}
